\documentclass[11pt,a4paper]{article}

\usepackage[T1]{fontenc}
\usepackage[utf8]{inputenc}
\usepackage{lmodern}
\usepackage[margin=1in]{geometry}
\usepackage{amsmath,amssymb,amsthm}
\usepackage{array}
\usepackage[hidelinks]{hyperref}

\newcommand{\smat}[4]{\left(\begin{smallmatrix}#1&#2\\#3&#4\end{smallmatrix}\right)}

\newtheorem{theorem}{Theorem}[section]
\newtheorem{proposition}[theorem]{Proposition}
\newtheorem{lemma}[theorem]{Lemma}
\newtheorem{corollary}[theorem]{Corollary}
\theoremstyle{definition}
\newtheorem{definition}[theorem]{Definition}
\newtheorem{remark}[theorem]{Remark}

\newcommand{\Cl}{\mathrm{Cl}}
\newcommand{\End}{\mathrm{End}}
\newcommand{\Aut}{\mathrm{Aut}}
\newcommand{\GL}{\mathrm{GL}}
\newcommand{\SL}{\mathrm{SL}}
\newcommand{\PGL}{\mathrm{PGL}}
\newcommand{\Sp}{\mathrm{Sp}}
\newcommand{\Orth}{\mathrm{O}}
\newcommand{\Aff}{\mathrm{Aff}}
\newcommand{\rad}{\mathrm{rad}}
\newcommand{\id}{\mathrm{id}}

\title{\textbf{Minimal Ideals, Spinor Dimensions, and Ideal Automorphism Groups in Clifford Algebras:\\[4pt]
\large Positive Split Regime and Complete Refutation of Conjecture 00000007628}}
\author{Formal Verification \& Mathematical Analysis Report}
\date{November 2026}

\begin{document}
\maketitle

\begin{abstract}
We investigate Conjecture \texttt{00000007628}, which asserts that for minimal ideals of Clifford algebras: (I)~the minimal-ideal dimension equals the spinor dimension, governed by an exact ideal-dimension conversion constant; (II)~this dimension law is exact and tested by semisimple decomposition; and (III)~the automorphism group of the minimal ideal is the symplectic group. We first establish a \emph{Positive Non-Vacuity Theorem} showing where this heuristic arises: on the nondegenerate split Clifford algebra $\Cl_{1,1}(K) \cong M_2(K)$ ($\mathrm{char}(K) \neq 2$), orthogonal primitive idempotents decompose $\Cl_{1,1}(K) = S_1 \oplus S_2$ into minimal left ideals of exact spinor dimension $2 = 2^{2/2}$ (conversion constant $4/2 = 2$), and $S_1$ carries a canonical Clifford-conjugation-covariant alternating bilinear form $\omega$ whose form-preserving linear automorphism group $\Aut_K(S_1, \omega)$ is literally $\Sp_2(K)$. Nevertheless, we prove that the unqualified conjecture is \textbf{false}: every one of Clauses (I)--(III) fails under standard mathematical definitions, and Clause (III) fails universally even on nondegenerate split Clifford algebras across all standard notions of the automorphism group of an ideal ($\Cl$-module automorphisms, $K$-linear automorphisms, non-unital left-ideal algebra automorphisms, unital two-sided-ideal algebra automorphisms, and canonical spinor-pairing isometries across mod-$8$ signatures). All core algebraic theorems and counterexamples are formally verified from first principles in Core Lean~4.19.0 with zero external dependencies and independently checked by an exact Python verification suite.
\end{abstract}

\section{Introduction and Taxonomic Analysis of the Conjecture}

Conjecture \texttt{00000007628} (\texttt{SOURCE.md}) states verbatim in English:
\begin{quote}
\emph{Definition: Minimal ideals with the ideal dimension law of Clifford algebras. Conjecture: The conversion constant is the ideal dimension law and the minimal-ideal dimension is exactly the spinor dimension; the law is exact with test by semisimple decomposition; and the automorphism group of the ideal is the symplectic group. (ideal-dimension conversion)}
\end{quote}
The companion Chinese formulation asserts the exact same three-part structure: defining minimal ideals with the Clifford ideal-dimension law (\emph{Clifford daishu de lixiang weishu lu}), it conjectures that (I)~the conversion constant is given by the ideal-dimension law and the minimal-ideal dimension equals the spinor dimension (\emph{jixiao lixiang weishu qia wei xuanliang weishu}); (II)~the law is exact and tested by semisimple decomposition (\emph{banshan fenjie}); and (III)~the automorphism group of the ideal is the symplectic group (\emph{lixiang de zitonggouqun wei xinqun}).

Let $K$ be a field and $(V, Q)$ an $n$-dimensional quadratic space over $K$. By the Poincar\'e--Birkhoff--Witt / Chevalley theorem, the Clifford algebra $\Cl(V, Q) = T(V)/\langle v \otimes v - Q(v)\cdot 1 \rangle$ has total dimension $\dim_K \Cl(V, Q) = 2^n$ for \emph{every} quadratic form $Q$, whether nondegenerate or degenerate. In classical geometric and physical spinor theory (Dirac and Weyl spinors over $\mathbb{C}$), the reference spinor dimension is $d_{\mathrm{spin}}(n) = 2^{\lfloor n/2 \rfloor}$; in Chevalley's and Lounesto's algebraic spinor theory~\cite{Chevalley1954,Lounesto2001}, spinors are elements of a minimal nonzero \emph{left} ideal $S = \Cl(V, Q)e$ generated by a primitive idempotent $e^2 = e$, and when $\Cl(V, Q)$ is semisimple of homogeneous simple-module dimension $\dim_K S$, the Wedderburn--Artin decomposition $\Cl(V, Q) \cong \bigoplus_{j=1}^{m_{\mathrm{id}}} S_j$ yields the exact conversion constant $m_{\mathrm{id}} = \dim_K \Cl(V, Q) / \dim_K S$.

To evaluate Conjecture \texttt{00000007628} rigorously without ambiguity, one must test every standard mathematical reading of its constituent terms:
\begin{enumerate}
  \item[\textbf{(R1)}] \textbf{Minimal two-sided vs.\ minimal left ideals:} In standard ring theory, an unqualified \emph{ideal} of $A = \Cl(V, Q)$ means a two-sided ideal ($A J A \subseteq J$), whereas algebraic spinor spaces are minimal \emph{left} ideals ($A S \subseteq S$).
  \item[\textbf{(R2)}] \textbf{Quadratic form degeneracy and spinor dimension:} Does the ideal-dimension law with semisimple decomposition hold for general quadratic forms $Q$ (including degenerate ones such as $Q = 0$), and does $\dim_K S = 2^{\lfloor n/2 \rfloor}$ hold even when $Q$ is nondegenerate and $\Cl(V, Q)$ is semisimple?
  \item[\textbf{(R3)}] \textbf{Automorphism group of the ideal:} Given a minimal ideal $I \subseteq \Cl(V, Q)$ without auxiliary extra data specified in the conjecture, there are five canonical mathematical definitions of ``the automorphism group of $I$'':
  \begin{enumerate}
    \item[(a)] \emph{Left $\Cl(V, Q)$-module automorphisms} $\Aut_{\Cl}(S) = \{ \varphi \in \GL_K(S) : \varphi(x\psi) = x\varphi(\psi)$ for all $x \in \Cl(V, Q)\}$;
    \item[(b)] \emph{$K$-linear automorphisms} $\GL_K(S)$ of the underlying vector space;
    \item[(c)] \emph{Non-unital $K$-algebra automorphisms} $\Aut_{K\text{-alg}}(S)$ preserving the induced subring multiplication $(\psi, \varphi) \mapsto \psi\varphi$ on a minimal left ideal $S$;
    \item[(d)] \emph{Unital $K$-algebra automorphisms} $\Aut_{K\text{-alg}}(J)$ of a minimal two-sided ideal $J$;
    \item[(e)] \emph{Spinor bilinear pairing automorphisms} $\Aut_K(S, \beta)$ preserving a canonical reversion- or Clifford-conjugation-covariant bilinear form $\beta$ on $S$~\cite{Porteous1995,Lounesto2001}.
  \end{enumerate}
\end{enumerate}

We prove that while a specialized combination---split signature $\Cl_{1,1}(K)$ with minimal \emph{left} ideals and pairing-isometry group $\Aut_K(S, \omega)$ under Reading~(R3e)---realizes all three clauses simultaneously (Section~\ref{sec:positive}), the unqualified conjecture fails across \textbf{all} general readings, and Clause~(III) fails under \textbf{every} one of (R3a)--(R3e) (Sections~\ref{sec:clauses12}--\ref{sec:automorphisms}).

\section{Positive Non-Vacuity Theorem on Split \texorpdfstring{$\Cl_{1,1}(K) \cong M_2(K)$}{Cl(1,1)(K)}}\label{sec:positive}

We first isolate the precise algebraic mechanism that motivates Conjecture \texttt{00000007628}.

\begin{theorem}[Positive Non-Vacuity in Split Signature $(1,1)$]\label{thm:positive}
Let $K$ be a field of characteristic $\neq 2$, and let $Q(x_1, x_2) = x_1^2 - x_2^2$ on $V = K^2$. Realize $\Cl_{1,1}(K) \cong M_2(K)$ via the matrix generators
\[
  \gamma_1 = \begin{pmatrix} 1 & 0 \\ 0 & -1 \end{pmatrix}, \qquad
  \gamma_2 = \begin{pmatrix} 0 & 1 \\ -1 & 0 \end{pmatrix}, \qquad
  \gamma_1^2 = I_2, \quad \gamma_2^2 = -I_2, \quad \gamma_1 \gamma_2 + \gamma_2 \gamma_1 = 0.
\]
\begin{enumerate}
  \item[(1)] \textbf{Semisimple decomposition and exact spinor dimension:} The elements $E_{11} = \frac{1}{2}(I_2 + \gamma_1)$ and $E_{22} = \frac{1}{2}(I_2 - \gamma_1)$ are orthogonal primitive idempotents ($E_{11}^2 = E_{11}$, $E_{22}^2 = E_{22}$, $E_{11}E_{22} = E_{22}E_{11} = 0$, $E_{11} + E_{22} = I_2$). They decompose $\Cl_{1,1}(K) = S_1 \oplus S_2$ into two isomorphic minimal nonzero left ideals $S_i = M_2(K)E_{ii}$ of dimension $\dim_K S_i = 2 = 2^{2/2}$, with exact conversion constant $\dim_K \Cl_{1,1}(K)/\dim_K S_1 = 4/2 = 2$.
  \item[(2)] \textbf{Canonical symplectic spinor pairing:} Under this isomorphism, Clifford conjugation $\overline{\,\cdot\,}$ (the unique anti-involution satisfying $\overline{\gamma_1} = -\gamma_1$ and $\overline{\gamma_2} = -\gamma_2$) is the matrix adjugate
  \[
    \overline{\begin{pmatrix} a & b \\ c & d \end{pmatrix}} = \begin{pmatrix} d & -b \\ -c & a \end{pmatrix} = \gamma_2^{-1} X^\top \gamma_2.
  \]
  Identifying $S_1 = \bigl\{\smat{u_1}{0}{u_2}{0}\bigr\} \cong K^2$, the bilinear form $\omega(u, v) = u_1 v_2 - u_2 v_1$ (coming from $u^\top \gamma_2 v = \omega(u, v) E_{11}$) is nondegenerate, alternating, and Clifford-conjugation covariant:
  \[
    \omega(X u, v) = \omega(u, \overline{X} v) \quad \forall X \in \Cl_{1,1}(K),\; u, v \in S_1.
  \]
  Moreover, $\omega(g u, g v) = (\det g)\,\omega(u, v)$ for all $g \in \End_K(S_1) \cong M_2(K)$, so the pairing-preserving linear automorphism group is literally the symplectic group:
  \[
    \Aut_K(S_1, \omega) = \{ g \in M_2(K) : \det g = 1 \} = \Sp_2(K).
  \]
\end{enumerate}
\end{theorem}

\begin{proof}
Direct matrix computation gives $\gamma_1^2 = I_2$, $\gamma_2^2 = -I_2$, and $\gamma_1\gamma_2 = \smat{0}{1}{1}{0} = -\gamma_2\gamma_1$. Since $\mathrm{char}(K) \neq 2$, $1, \gamma_1, \gamma_2, \gamma_1\gamma_2$ span $M_2(K)$, confirming $\Cl_{1,1}(K) \cong M_2(K)$ and $E_{11} = \frac{1}{2}(I_2 + \gamma_1) = \smat{1}{0}{0}{0}$, $E_{22} = \frac{1}{2}(I_2 - \gamma_1) = \smat{0}{0}{0}{1}$. Every $X \in M_2(K)$ splits uniquely as $X = X E_{11} + X E_{22} \in S_1 \oplus S_2$. For any $0 \neq u = (u_1, u_2)^\top \in S_1$ and any $w = (w_1, w_2)^\top \in S_1$, if $u_1 \neq 0$ then $\smat{w_1 u_1^{-1}}{0}{w_2 u_1^{-1}}{0} u = w$, while if $u_1 = 0$ then $u_2 \neq 0$ and $\smat{0}{w_1 u_2^{-1}}{0}{w_2 u_2^{-1}} u = w$; thus $S_1$ is simple of dimension $2$.

For (2), $\overline{X} = \gamma_2^{-1} X^\top \gamma_2$ is an anti-involution fixing $I_2$ and negating $\gamma_1, \gamma_2, \gamma_1\gamma_2$. Expanding $\omega(Xu, v)$ and $\omega(u, \overline{X}v)$ gives $(a u_1 + b u_2)v_2 - (c u_1 + d u_2)v_1 = u_1(-c v_1 + a v_2) - u_2(d v_1 - b v_2)$. Applying covariance with $\overline{g}\,g = (\det g)I_2$ yields $\omega(gu, gv) = \omega(u, \overline{g}gv) = (\det g)\,\omega(u, v)$. Evaluating on the standard symplectic basis $e_1 = (1,0)^\top, e_2 = (0,1)^\top$ with $\omega(e_1, e_2) = 1$ shows that $g$ preserves $\omega$ if and only if $\det g = 1$.
\end{proof}

\begin{remark}[Compact Quaternionic Norm-$1$ Subgroup on $\Cl_{0,2}(\mathbb{R}) \cong \mathbb{H}$]
A second classical source of symplectic terminology in low-dimensional spinor theory occurs on $\Cl_{0,2}(\mathbb{R}) \cong \mathbb{H}$, where $\End_{\Cl}(\mathbb{H}) \cong \mathbb{H}^{\mathrm{op}}$ has unit group $\mathbb{H}^\times \cong \mathbb{R}_{>0} \times \Sp(1)$; restricting to unit-norm quaternions yields the compact symplectic group $\Sp(1) \cong \mathrm{SU}(2)$, though the full module automorphism group $\mathbb{H}^\times$ has dimension $4 \neq 3 = \dim_{\mathbb{R}} \Sp(1)$ and $\dim_{\mathbb{R}} \mathbb{H} = 4 \neq 2 = 2^{2/2}$.
\end{remark}

\section{Refutation of Clauses I and II: Minimal-Ideal Dimension and Semisimplicity}\label{sec:clauses12}

We now show that Clauses~(I) and~(II) fail whenever ``minimal ideal'' is read as minimal two-sided ideal, whenever $Q$ is degenerate, or whenever $\Cl(V, Q)$ has non-split division-algebra Wedderburn components.

\begin{proposition}[Failure for Minimal Two-Sided Ideals]\label{prop:twosided}
Over any field $K$, $M_2(K)$ is a simple two-sided $K$-algebra: for every nonzero $X \in M_2(K)$, there exist $L_1, R_1, L_2, R_2 \in M_2(K)$ such that $L_1 X R_1 + L_2 X R_2 = I_2$. Consequently, for $\mathrm{char}(K) \neq 2$, the unique minimal nonzero two-sided ideal of $\Cl_{1,1}(K) \cong M_2(K)$ is $\Cl_{1,1}(K)$ itself, of dimension
\[
  \dim_K J_{\min} = 4 \neq 2 = d_{\mathrm{spin}}(2).
\]
More generally, for $\Cl_{m,m}(K) \cong M_{2^m}(K)$ ($m \ge 1$), the unique minimal nonzero two-sided ideal has dimension $4^m = (d_{\mathrm{spin}})^2 > 2^m = d_{\mathrm{spin}}$.
\end{proposition}

\begin{proof}
If $X_{ij} \neq 0$ for some $i, j \in \{1, 2\}$, then $X_{ij}^{-1} E_{1i} X E_{j1} + X_{ij}^{-1} E_{2i} X E_{j2} = E_{11} + E_{22} = I_2$. Thus every nonzero two-sided ideal contains $I_2$ and equals $M_2(K)$ (and similarly $M_{2^m}(K)$ for all $m \ge 1$).
\end{proof}

\begin{theorem}[Failure of Semisimplicity and Spinor Dimension for Degenerate $Q$]\label{thm:exterior}
Let $K$ be any field and $n \ge 2$, and consider the zero quadratic form $Q \equiv 0$ on $V = K^n$, whose Clifford algebra is the exterior algebra $\Cl(K^n, 0) = \Lambda(K^n)$ of dimension $2^n$.
\begin{enumerate}
  \item[(1)] \textbf{Failure of semisimplicity:} The augmentation ideal $\Lambda^+(K^n) = \bigoplus_{k=1}^n \Lambda^k(K^n)$ is a nilpotent two-sided ideal of dimension $2^n - 1 > 0$ (for $n = 2$, $z^2 = 0$ for all $z \in \Lambda^+(K^2)$). The only idempotents in $\Lambda(K^n)$ are $0$ and $1$, and the top exterior power $I_{\mathrm{top}} = \Lambda^n(K^n)$ admits \emph{no} left-ideal direct-sum complement $L$ in $\Lambda(K^n)$.
  \item[(2)] \textbf{Failure of spinor dimension even for minimal left ideals:} Every nonzero left ideal of $\Lambda(K^n)$ contains $\Lambda^n(K^n)$. Hence $I_{\mathrm{top}} = \Lambda^n(K^n)$ is the \emph{unique} minimal nonzero left ideal (and unique minimal nonzero two-sided ideal) of $\Lambda(K^n)$, with dimension
  \[
    \dim_K I_{\mathrm{top}} = 1 \neq 2^{\lfloor n/2 \rfloor} \qquad (\text{e.g., } 1 \neq 2 \text{ at } n = 2).
  \]
\end{enumerate}
\end{theorem}

\begin{proof}
Let $\{e_1, \dots, e_n\}$ be a basis of $K^n$ and $e_{\mathrm{top}} = e_1 \wedge \cdots \wedge e_n$. Given any nonzero element $x = \sum_S c_S e_S \in \Lambda(K^n)$ indexed by $S \subseteq \{1,\dots,n\}$, pick $S_0$ of minimal cardinality $|S_0|$ with $c_{S_0} \neq 0$, and set $S_0^{\mathrm{c}} = \{1,\dots,n\} \setminus S_0$. For any $S \neq S_0$ with $c_S \neq 0$, we have $|S| \ge |S_0|$ and $S \neq S_0$, so $S \cap S_0^{\mathrm{c}} \neq \emptyset$ and therefore $e_{S_0^{\mathrm{c}}} \wedge e_S = 0$, whereas $e_{S_0^{\mathrm{c}}} \wedge e_{S_0} = \pm e_{\mathrm{top}}$. Hence $\pm c_{S_0}^{-1} e_{S_0^{\mathrm{c}}} \wedge x = e_{\mathrm{top}}$, proving (2).

For (1), if $a = a_0 + z \in \Lambda(K^n)$ ($a_0 \in K$, $z \in \Lambda^+$) satisfies $a^2 = a$, degree-$0$ comparison gives $a_0^2 = a_0$, so $a_0 \in \{0, 1\}$. If $a_0 = 0$, $a = a^{2^k} = 0$ by nilpotency of $\Lambda^+$; if $a_0 = 1$, write $z = z_r + z_{>r}$ with lowest nonzero homogeneous piece $z_r \in \Lambda^r$ ($r \ge 1$): comparing degree $r$ in $(1 + z)^2 = 1 + z$ gives $z_r + z_r = z_r \implies z_r = 0$, so $z = 0$ (valid in \emph{every} characteristic!). Finally, if $\Lambda(K^n) = I_{\mathrm{top}} \oplus L$ for some left ideal $L$, write $1 = u + \ell$ with $u \in I_{\mathrm{top}}$ and $\ell \in L$. Since $e_{\mathrm{top}} \wedge u = 0$, left-multiplying by $e_{\mathrm{top}}$ gives $e_{\mathrm{top}} = e_{\mathrm{top}} \wedge 1 = e_{\mathrm{top}} \wedge \ell \in I_{\mathrm{top}} \cap L = \{0\}$, a contradiction.
\end{proof}

\begin{proposition}[Failure of $2^{\lfloor n/2 \rfloor}$ on Nondegenerate Semisimple Division Clifford Algebras]\label{prop:quat}
For the nondegenerate negative-definite form $Q(x_1, x_2) = -x_1^2 - x_2^2$ over $K = \mathbb{Q}$ (or $K = \mathbb{R}$), $\Cl_{0,2}(K) \cong \mathbb{H}_K$ is a semisimple $4$-dimensional division algebra: every nonzero $q = a + bi + cj + dk \in \mathbb{H}_K$ satisfies $\overline{q}\,q = a^2 + b^2 + c^2 + d^2 > 0$, hence is invertible with $q^{-1} = (a^2 + b^2 + c^2 + d^2)^{-1}\overline{q}$. Therefore, the unique minimal nonzero left ideal of $\Cl_{0,2}(K)$ is $\mathbb{H}_K$ itself, of dimension $\dim_K \mathbb{H}_K = 4 \neq 2 = 2^{\lfloor 2/2 \rfloor}$ (and similarly $\Cl_{0,1}(\mathbb{R}) \cong \mathbb{C}$ has $\dim_{\mathbb{R}} \mathbb{C} = 2 \neq 1 = 2^{\lfloor 1/2 \rfloor}$).
\end{proposition}

\section{Refutation of Clause III: Ideal Automorphism Groups}\label{sec:automorphisms}

Even if one restricts to nondegenerate split Clifford algebras $\Cl_{m,m}(K) \cong M_{2^m}(K)$ ($\mathrm{char}(K) \neq 2$) and minimal \emph{left} ideals $S$---where Clauses~(I) and~(II) hold with $\dim_K S = 2^m$---Clause~(III) fails across \textbf{all five} standard interpretations (R3a)--(R3e).

\subsection{(R3a) Left \texorpdfstring{$\Cl$}{Cl}-Module Automorphisms \texorpdfstring{$\Aut_{\Cl}(S)$}{Aut\_Cl(S)}}

\begin{theorem}[Non-Isomorphism of $\Aut_{\Cl_{1,1}(K)}(S_1)$ with Any Symplectic Group]\label{thm:clmod}
Let $K$ and $F$ be arbitrary fields, and let $S_1 = M_2(K)E_{11} \cong K^2$ be the minimal left ideal of $M_2(K)$ ($\cong \Cl_{1,1}(K)$ when $\mathrm{char}(K) \neq 2$).
\begin{enumerate}
  \item[(1)] Every map $\varphi : S_1 \to S_1$ satisfying $\varphi(X u) = X \varphi(u)$ for all $X \in M_2(K)$ and $u \in S_1$ is scalar multiplication $\varphi(u) = c u$ with $c = (\varphi(e_1))_1 \in K$. Thus $\Aut_{M_2(K)}(S_1) \cong K^\times$ is \textbf{abelian}.
  \item[(2)] For every field $F$ and every rank $r \ge 1$, $\Sp_{2r}(F)$ is \textbf{non-abelian}: in $\Sp_2(F)$ ($\hookrightarrow \Sp_{2r}(F)$), the transvections $U_{12} = \smat{1}{1}{0}{1}$ and $L_{21} = \smat{1}{0}{1}{1}$ satisfy $\det U_{12} = \det L_{21} = 1$, yet $(U_{12} L_{21})_{11} = 1 + 1 \neq 1 = (L_{21} U_{12})_{11}$ since $1 \neq 0$ in $F$.
  \item[(3)] Consequently, $\Aut_{\Cl_{1,1}(K)}(S_1)$ is \textbf{never} isomorphic as a group to $\Sp_{2r}(F)$ for any fields $K, F$ and any $r \ge 1$.
\end{enumerate}
\end{theorem}

\begin{proof}
Since $E_{11} e_1 = e_1$, equivariance gives $\varphi(e_1) = E_{11}\varphi(e_1) = (c, 0)^\top$ where $c = (\varphi(e_1))_1$. For any $u = (u_1, u_2)^\top \in S_1$, let $X_u = \smat{u_1}{0}{u_2}{0}$; then $X_u e_1 = u$, so $\varphi(u) = X_u \varphi(e_1) = (u_1 c, u_2 c)^\top = c u$. Thus any two module automorphisms commute, whereas $U_{12}, L_{21} \in \Sp_2(F)$ do not commute; since commutativity is preserved under group isomorphism, no isomorphism exists.
\end{proof}

\begin{remark}[Bott Periodicity Over $\mathbb{R}$ and Radial Dilations]
Over $\mathbb{R}$, Bott periodicity restricts $\End_{\Cl_{p,q}(\mathbb{R})}(S)$ to $\mathbb{R}, \mathbb{C}$, or $\mathbb{H}$, so $\Aut_{\Cl_{p,q}(\mathbb{R})}(S) \in \{\mathbb{R}^\times, \mathbb{C}^\times, \mathbb{H}^\times\}$ has Lie dimension $1, 2$, or $4$ for \emph{all} $(p, q)$, whereas $\dim_{\mathbb{R}} \Sp_{2r}(\mathbb{R}) = \dim_{\mathbb{R}} \Sp(r) = r(2r+1) \in \{3, 10, 21, 36, \dots\}$ never equals $1, 2$, or $4$. Moreover, for any field $K$ with $|K| > 3$, radial dilations $\lambda \cdot \id_S \in \Aut_{\Cl}(S)$ ($\lambda \neq \pm 1$) satisfy $\beta(\lambda x, \lambda y) = \lambda^2 \beta(x, y) \neq \beta(x, y)$ whenever $\beta(x, y) \neq 0$, so $\Aut_{\Cl}(S)$ preserves \emph{no} nonzero bilinear form on $S$.
\end{remark}

\subsection{(R3b) \texorpdfstring{$K$}{K}-Linear Automorphisms \texorpdfstring{$\GL_K(S)$}{GL\_K(S)}}

\begin{proposition}[Failure for $\GL_K(S)$]\label{prop:gl}
Let $K$ be a field with $\mathrm{char}(K) \neq 2$.
\begin{enumerate}
  \item[(1)] On $\Cl_{1,0}(K) = K[e_1]/(e_1^2 - 1) \cong K \oplus K$ ($n = 1$, nondegenerate $Q(x) = x^2$), the minimal ideals $I_\pm = K \cdot \frac{1}{2}(1 \pm e_1)$ have odd dimension $\dim_K I_\pm = 1 = 2^{\lfloor 1/2 \rfloor}$. Every alternating bilinear form $\omega : I_+ \times I_+ \to K$ satisfies $\omega(x, y) = xy\,\omega(1, 1) = 0$, so $I_+$ admits \textbf{no} symplectic structure, and $\GL_K(I_+) \cong K^\times$ is abelian.
  \item[(2)] On $\Cl_{1,1}(K) \cong M_2(K)$ ($\dim_K S_1 = 2$), $\mathrm{diag}(-1, 1) \in \GL_2(K)$ has $\det = -1 \neq 1$, so $\GL_2(K) \supsetneq \Sp_2(K)$ strictly; over $K = \mathbb{F}_3$, $|\GL_2(\mathbb{F}_3)| = 48 \neq 24 = |\Sp_2(\mathbb{F}_3)|$. For $N = 2^m \ge 4$ ($\Cl_{m,m}(K)$), $\dim \GL_N = N^2 > \frac{N(N+1)}{2} = \dim \Sp_N$ (e.g., $16 > 10$ at $N = 4$).
\end{enumerate}
\end{proposition}

\subsection{(R3c) Non-Unital \texorpdfstring{$K$}{K}-Algebra Automorphisms \texorpdfstring{$\Aut_{K\text{-alg}}(S)$}{Aut\_(K-alg)(S)} of a Minimal Left Ideal}

Because a minimal left ideal $S = M_N(K)E_{11}$ ($N = 2^m \ge 2$) is closed under matrix multiplication, it forms a non-unital $K$-algebra with product $u \cdot v = v_1 u$ for $u, v \in K^N$.

\begin{theorem}[Structure and Non-Symplecticity of $\Aut_{K\text{-alg}}(S)$]\label{thm:s1alg}
Let $K$ be any field and $S = M_N(K)E_{11} \cong K^N$ ($N = 2^m \ge 2$).
\begin{enumerate}
  \item[(1)] A linear bijection $g \in \GL_N(K)$ satisfies $g(u \cdot v) = g(u) \cdot g(v)$ for all $u, v \in S$ if and only if $g$ has block form $\smat{1}{0}{q}{B}$ with $q \in K^{N-1}$ and $B \in \GL_{N-1}(K)$. Thus $\Aut_{K\text{-alg}}(S) \cong K^{N-1} \rtimes \GL_{N-1}(K)$ has algebraic dimension $N(N - 1)$, which \textbf{never} equals $\dim \Sp_N = \frac{N(N+1)}{2}$ for any even integer $N = 2r \ge 2$, since
  \[
    2r(2r - 1) - r(2r + 1) = r(2r - 3) \neq 0 \quad \forall r \in \mathbb{Z}_{\ge 1}.
  \]
  \item[(2)] At $N = 2$ ($S_1 \subset \Cl_{1,1}(K)$), $\Aut_{K\text{-alg}}(S_1) \cong \Aff_1(K) = K \rtimes K^\times$ is \textbf{metabelian} ($2$-step solvable: every commutator has diagonal part $1$, and elements with diagonal part $1$ commute). Conversely, for \textbf{every} field $F$ with $\mathrm{char}(F) \neq 2$, the derived subgroup $[\Sp_2(F), \Sp_2(F)]$ is \textbf{non-abelian}: setting $J_0 = \smat{0}{-1}{1}{0} \in \Sp_2(F)$, the commutators
  \[
    X = [J_0, U_{12}] = \begin{pmatrix} 1 & -1 \\ -1 & 2 \end{pmatrix}, \qquad
    Y = [J_0, L_{21}] = \begin{pmatrix} 2 & -1 \\ -1 & 1 \end{pmatrix}
  \]
  satisfy $(XY)_{12} = -2 \neq -4 = (YX)_{12}$ since $2 \neq 0$ in $F$. Hence $\Aut_{K\text{-alg}}(S_1) \not\cong \Sp_2(F)$ for all fields $K, F$ with $\mathrm{char}(F) \neq 2$ (and over $\mathbb{F}_3$, $|\Aff_1(\mathbb{F}_3)| = 6 \neq 24 = |\Sp_2(\mathbb{F}_3)|$).
\end{enumerate}
\end{theorem}

\begin{proof}
In $S \cong K^N$, $\rad(S) = \{u : u_1 = 0\} = \{u : u \cdot u = 0\}$ is characteristic of dimension $N - 1$, forcing the first row of $g$ to vanish outside entry $(1,1)$; applying $g(e_1) = g(e_1 \cdot e_1) = g(e_1)_1 g(e_1)$ with $g(e_1)_1 \neq 0$ yields $g_{11} = 1$, and any $\smat{1}{0}{q}{B}$ preserves $u \cdot v = v_1 u$. Part~(2) follows by direct matrix multiplication and preservation of derived-subgroup commutativity under group isomorphism.
\end{proof}

\subsection{(R3d) Unital \texorpdfstring{$K$}{K}-Algebra Automorphisms \texorpdfstring{$\Aut_{K\text{-alg}}(M_N(K))$}{Aut\_(K-alg)(M\_N(K))} of a Minimal Two-Sided Ideal}

\begin{theorem}[Center Obstruction for $\Aut_{K\text{-alg}}(M_2(K))$]\label{thm:twosided_aut}
For \textbf{every} field $K$, the algebra automorphism group $\Aut_{K\text{-alg}}(M_2(K))$ ($\cong \PGL_2(K)$) has trivial center $\{\id\}$. Conversely, for every field $F$ of characteristic $\neq 2$ and every $r \ge 1$, $-I_{2r} \in Z(\Sp_{2r}(F))$ satisfies $-I_{2r} \neq I_{2r}$. Hence $\Aut_{K\text{-alg}}(M_2(K)) \not\cong \Sp_{2r}(F)$ for all fields $K$, all fields $F$ with $\mathrm{char}(F) \neq 2$, and all $r \ge 1$.
\end{theorem}

\begin{proof}
We give a self-contained proof without invoking Skolem--Noether. Suppose $\alpha$ lies in the center of $\Aut_{K\text{-alg}}(M_2(K))$. For any invertible $g \in \GL_2(K)$, $\alpha$ commutes with $\mathrm{Ad}_g(Y) = g Y g^{-1}$. Setting $Y = \alpha^{-1}(X)$ gives $\alpha(g) X \alpha(g^{-1}) = g X g^{-1}$, whence $M = g^{-1}\alpha(g)$ commutes with all $X \in M_2(K)$. Commuting with $E_{12}$ and $E_{21}$ forces $M = c I_2$, so $\alpha(g) = c g$. Applying this to $g_1 = I_2 + E_{12}$ yields $\alpha(E_{12}) = \alpha(g_1 - I_2) = (c_1 - 1)I_2 + c_1 E_{12}$. Since $E_{12}^2 = 0$, $0 = \alpha(E_{12})^2$ has $(1,1)$-entry $(c_1 - 1)^2 = 0$, so $c_1 = 1$ and $\alpha(E_{12}) = E_{12}$. Identically, $g_2 = I_2 + E_{21}$ gives $\alpha(E_{21}) = E_{21}$, whence $\alpha(E_{11}) = \alpha(E_{12}E_{21}) = E_{11}$ and $\alpha(E_{22}) = \alpha(E_{21}E_{12}) = E_{22}$. By $K$-linearity, $\alpha = \id$. Since group isomorphisms preserve centers, $\Aut_{K\text{-alg}}(M_2(K))$ cannot be isomorphic to $\Sp_{2r}(F)$ (noting that over $\mathbb{F}_3$ and $\mathbb{F}_5$, $|\PGL_2(\mathbb{F}_q)| = |\Sp_2(\mathbb{F}_q)| \in \{24, 120\}$, yet their centers have orders $1 \neq 2$). Moreover, for $N = 2^m \ge 4$, $\dim \PGL_N = N^2 - 1 \neq \frac{N(N+1)}{2} = \dim \Sp_N$ ($15 \neq 10$ at $N = 4$).
\end{proof}

\subsection{(R3e) Spinor Bilinear Invariants and Mod-8 Periodicity}

Finally, even if one enriches a minimal left ideal $S$ with its canonical $\mathrm{Pin}$-covariant spinor bilinear form $\beta(\psi, \varphi)$ satisfying $\beta(X\psi, \varphi) = \beta(\psi, \tau(X)\varphi)$ for $\tau \in \{\text{reversion } \widetilde{\cdot},\, \text{Clifford conjugation } \overline{\cdot}\}$~\cite{Porteous1995,Lounesto2001}, the resulting isometry group $\Aut_K(S, \beta)$ is \emph{not} always symplectic.

\begin{theorem}[Euclidean $\Cl_{2,0}(K)$ Reversion Orthogonality and Mod-$8$ Double Orthogonality]\label{thm:mod8}
Let $K$ be a field with $\mathrm{char}(K) \neq 2$.
\begin{enumerate}
  \item[(1)] \textbf{Symmetric reversion form on $\Cl_{2,0}(K) \cong M_2(K)$:} Realize $\Cl_{2,0}(K)$ ($Q(x_1, x_2) = x_1^2 + x_2^2$) via $\varepsilon_1 = \smat{1}{0}{0}{-1}$ and $\varepsilon_2 = \smat{0}{1}{1}{0}$. Reversion is matrix transpose $\widetilde{X} = X^\top$, and \emph{every} reversion-covariant bilinear form $\beta$ on $S_1 \cong K^2$ has the symmetric form $\beta(u, v) = c(u_1 v_1 + u_2 v_2)$ ($c = \beta(e_1, e_1) \in K$). When $c \neq 0$, $\beta$ is never alternating and its isometry group is the orthogonal group $\Orth_2(K)$ (of order $|\Orth_2(\mathbb{F}_3)| = 8 \neq 24 = |\Sp_2(\mathbb{F}_3)|$, and compact $1$-dimensional $\Orth_2(\mathbb{R})$ vs.\ $3$-dimensional $\Sp_2(\mathbb{R})$).
  \item[(2)] \textbf{Double orthogonality at $n \equiv 0 \pmod{8}$:} On $\Cl_{m,m}(K) \cong M_N(K)$ ($n = 2m$, $N = 2^m$), any anti-involution $\tau(X) = B^{-1} X^\top B$ ($B^\top = \epsilon B$, $\epsilon \in \{+1, -1\}$) has $+1$-eigenspace of dimension $\frac{N(N+1)}{2}$ if $\epsilon = +1$ ($\Orth_N$) and $\frac{N(N-1)}{2}$ if $\epsilon = -1$ ($\Sp_N$). Counting $k$-blades fixed by reversion ($k \equiv 0, 1 \pmod 4$) and Clifford conjugation ($k \equiv 0, 3 \pmod 4$) gives
  \[
    \epsilon_{\mathrm{rev}}(2m) = (-1)^{m(m-1)/2}, \qquad \epsilon_{\mathrm{conj}}(2m) = (-1)^{m(m+1)/2}.
  \]
  Whenever $m \equiv 0 \pmod{4}$---in particular at $n = 8$ ($\Cl_{4,4}(K) \cong M_{16}(K)$, $N = 16$), where both anti-involutions fix $136 = \frac{16 \cdot 17}{2} \neq 120 = \frac{16 \cdot 15}{2}$ basis blades---\textbf{both} canonical anti-involutions induce symmetric ($\Orth_N$) forms and \textbf{neither} induces a symplectic ($\Sp_N$) form.
\end{enumerate}
\end{theorem}

\begin{proof}
For (1), $\beta(E_{11} e_1, e_2) = \beta(e_1, E_{11}^\top e_2) = \beta(e_1, 0) = 0$ and similarly $\beta(e_2, e_1) = 0$, while $\beta(e_2, e_2) = \beta(\varepsilon_2 e_1, e_2) = \beta(e_1, \varepsilon_2^\top e_2) = \beta(e_1, e_1) = c$, giving $\beta(u, v) = c(u_1 v_1 + u_2 v_2)$. If $\beta(u, u) = 0$ for all $u$, then $c = \beta(e_1, e_1) = 0$, so $\beta \equiv 0$. For (2), $(-1)^{k(k-1)/2} = \mathrm{Re}((1-i)i^k)$ and $(-1)^{k(k+1)/2} = \mathrm{Re}((1+i)i^k)$, so summing against $\binom{2m}{k}$ gives $\mathrm{Tr}(\tau) = \mathrm{Re}((1 \mp i)(1+i)^{2m}) = 2^m \mathrm{Re}((1 \mp i)i^m) = \epsilon_\tau(2m)\,N$, matching $\frac{N(N+\epsilon_\tau)}{2} - \frac{N(N-\epsilon_\tau)}{2} = \epsilon_\tau N$. At $m = 4$ ($N = 16$), $i^4 = 1$ gives $\epsilon_{\mathrm{rev}}(8) = \epsilon_{\mathrm{conj}}(8) = +1$ and $\sum_{k \in \{0,1,4,5,8\}} \binom{8}{k} = 1 + 8 + 70 + 56 + 1 = 136 = \sum_{k \in \{0,3,4,7,8\}} \binom{8}{k}$.
\end{proof}

\begin{table}[ht]
\centering
\small
\renewcommand{\arraystretch}{1.2}
\caption{Complete taxonomy of readings of Conjecture \texttt{00000007628} and sharp counterexamples.}
\vspace{4pt}
\label{tab:taxonomy}
\begin{tabular}{@{}p{0.26\textwidth}p{0.23\textwidth}p{0.45\textwidth}@{}}
\hline
\textbf{Reading / Clause} & \textbf{Counterexample} & \textbf{Exact Failure Mechanism} \\
\hline
Minimal two-sided ideal $J$ & $\Cl_{1,1}(K) \cong M_2(K)$ & Simple algebra: $\dim_K J = 4 \neq 2 = d_{\mathrm{spin}}$ ($4^m \neq 2^m$). \\
Degenerate $Q \equiv 0$ & $\Cl(K^n, 0) = \Lambda(K^n)$, $n \ge 2$ & $\rad(\Lambda) = \Lambda^+ \neq 0$ (not semisimple); $\dim_K I_{\min} = 1 \neq 2^{\lfloor n/2 \rfloor}$. \\
Division $\Cl(V, Q)$ ($d_{\mathrm{spin}} = 2^{\lfloor n/2 \rfloor}$) & $\Cl_{0,2}(\mathbb{Q}) \cong \mathbb{H}_{\mathbb{Q}}$ & Division ring: $\dim_{\mathbb{Q}} I_{\min} = 4 \neq 2 = 2^{2/2}$. \\
(R3a) $\Aut_{\Cl}(S)$ & $\Cl_{1,1}(K)$, $S_1 \cong K^2$ & $\Aut_{\Cl}(S_1) \cong K^\times$ abelian vs.\ non-abelian $\Sp_{2r}(F)$ ($2 \neq 24$ on $\mathbb{F}_3$). \\
(R3b) $\GL_K(S)$ & $\Cl_{1,0}(K)$; $\Cl_{1,1}(\mathbb{F}_3)$ & Dim $1$ has no symplectic form; $|\GL_2(\mathbb{F}_3)| = 48 \neq 24$; $N^2 > \frac{N(N+1)}{2}$. \\
(R3c) $\Aut_{K\text{-alg}}(S)$ (left ideal) & $\Cl_{1,1}(K)$, $S_1 \cong K^2$ & $\Aff_1(K)$ metabelian ($|\Aff_1(\mathbb{F}_3)| = 6$) vs.\ non-metabelian $\Sp_2(F)$. \\
(R3d) $\Aut_{K\text{-alg}}(J)$ ($2$-sided) & $\Cl_{1,1}(K) \cong M_2(K)$ & $Z(\PGL_2(K)) = \{1\}$ vs.\ $\{\pm I\} \subseteq Z(\Sp_{2r}(F))$; $\dim \PGL_4 = 15 \neq 10$. \\
(R3e) Spinor pairing $\Aut_K(S, \beta)$ & $\Cl_{2,0}(K)$; $\Cl_{4,4}(K)$ ($n=8$) & $\Cl_{2,0}$ reversion gives $\Orth_2$ ($8 \neq 24$ on $\mathbb{F}_3$); at $n=8$ both give $\Orth_{16}$ ($136 \neq 120$). \\
\hline
\end{tabular}
\end{table}

\begin{table}[ht]
\centering
\footnotesize
\setlength{\tabcolsep}{4.5pt}
\renewcommand{\arraystretch}{1.2}
\caption{Exact finite-group order spectrum on $S_1 \cong \mathbb{F}_3^2$, algebraic dimensions on $S \cong K^N$ ($N = 2^m$), and Lounesto--Porteous mod-$8$ spinor form types on $\Cl_{m,m}(K) \cong M_{2^m}(K)$.}
\vspace{4pt}
\label{tab:spectrum}
\begin{tabular}{@{}lcclcc@{}}
\hline
\textbf{Group on $S_1 \cong \mathbb{F}_3^2$} & \textbf{$|G|$ ($\mathbb{F}_3$)} & \textbf{Dim ($N$)} & \textbf{$m \bmod 4$ ($n=2m$)} & \textbf{Reversion} & \textbf{Conjugation} \\
\hline
$\Aut_{\Cl_{1,1}}(S_1) \cong K^\times$ & $2$ & $1$ & $m \equiv 0$ ($n=0,8$) & $\Orth_N$ ($\epsilon=+1$) & $\Orth_N$ ($\epsilon=+1$) \\
$\Aut_{\text{alg}}(S_1) \cong \Aff_1$ & $6$ & $N(N-1)$ & $m \equiv 1$ ($n=2,10$) & $\Orth_N$ ($\epsilon=+1$) & $\Sp_N$ ($\epsilon=-1$) \\
$\Aut(S_1, \beta_{\mathrm{rev}}^{\Cl_{2,0}}) = \Orth_2$ & $8$ & $N(N-1)/2$ & $m \equiv 2$ ($n=4,12$) & $\Sp_N$ ($\epsilon=-1$) & $\Sp_N$ ($\epsilon=-1$) \\
$\Sp_2(\mathbb{F}_3) = \SL_2(\mathbb{F}_3)$ & $24$ & $N(N+1)/2$ & $m \equiv 3$ ($n=6,14$) & $\Sp_N$ ($\epsilon=-1$) & $\Orth_N$ ($\epsilon=+1$) \\
$\GL_2(\mathbb{F}_3)$ & $48$ & $N^2$ & \multicolumn{3}{l}{At $n=8, N=16$: both fix $136 = \frac{16\cdot 17}{2} \neq 120$ blades.} \\
\hline
\end{tabular}
\end{table}

\section{Formal Verification in Core Lean 4.19.0 and Computational Audit}\label{sec:verification}

All algebraic foundations, positive split-signature theorems, and multi-reading refutations are formalized in \textbf{Core Lean~4.19.0} (toolchain \nolinkurl{leanprover/lean4:v4.19.0}, strict compiler flag \nolinkurl{-DwarningAsError=true}, zero external dependencies, zero \texttt{sorry}, \texttt{admit}, or \nolinkurl{native_decide}, and zero custom axioms), as summarized in Table~\ref{tab:lean}.

\begin{table}[ht]
\centering
\footnotesize
\renewcommand{\arraystretch}{1.2}
\caption{Mapping between mathematical results and Core Lean~4.19.0 declarations.}
\vspace{4pt}
\label{tab:lean}
\begin{tabular}{@{}>{\raggedright\arraybackslash}p{0.34\textwidth}>{\raggedright\arraybackslash}p{0.60\textwidth}@{}}
\hline
\textbf{Lean Module \& Paper Result} & \textbf{Key Formalized Declarations} \\
\hline
\texttt{Field.lean} (Foundations) & \texttt{CommField}, \texttt{F3}, \texttt{MulEquiv}, \texttt{no\_mulEquiv\_of\_abelian\_nonabelian}, \texttt{no\_mulEquiv\_of\_trivial\_center}, \texttt{no\_bij\_of\_card\_ne} \\
\texttt{Mat2.lean} (Thm.~\ref{thm:positive}, Prop.~\ref{prop:twosided}, Thm.~\ref{thm:mod8}(1)) & \texttt{cl11\_semisimple\_decomposition}, \texttt{omega\_clifford\_covariant}, \texttt{preserves\_omega\_iff\_det\_eq\_one}, \texttt{mat2\_is\_simple\_two\_sided}, \texttt{cl20\_reversion\_form\_is\_symmetric} \\
\texttt{ExteriorAndQuat.lean} (Thm.~\ref{thm:exterior}, Prop.~\ref{prop:quat}, Prop.~\ref{prop:gl}(1)) & \texttt{idempotent\_trivial}, \texttt{e12\_mem\_every\_nonzero\_left\_ideal}, \texttt{ext2\_no\_semisimple\_complement}, \texttt{conj\_mul\_self\_and\_norm\_pos}, \texttt{one\_dim\_no\_alternating\_form} \\
\texttt{Automorphisms.lean} (Thms.~\ref{thm:clmod}--\ref{thm:mod8}) & \texttt{no\_iso\_clModAut\_sp2}, \texttt{gl2\_not\_equal\_sp2}, \texttt{s1\_alg\_aut\_iff\_affine}, \texttt{no\_iso\_s1AlgAut\_sp2}, \texttt{Mat2AlgAut.center\_trivial}, \texttt{no\_iso\_mat2AlgAut\_sp2}, \texttt{f3\_exact\_group\_orders}, \texttt{f3\_all\_four\_not\_bijective\_to\_sp2}, \texttt{cl8\_both\_involutions\_orthogonal\_not\_symplectic} \\
\texttt{Main.lean} (Synthesis) & \texttt{conjecture\_00000007628\_positive\_split\_case}, \texttt{conjecture\_00000007628\_complete\_refutation}, \texttt{conjecture\_00000007628\_refuted\_over\_F3} \\
\hline
\end{tabular}
\end{table}
In parallel, the standalone script \texttt{verify.py} independently checks all six verification suites using exact rational (\texttt{fractions.Fraction}) and finite-field ($\mathbb{F}_3, \mathbb{F}_5$) arithmetic. Conjecture \texttt{00000007628} is therefore definitively \textbf{disproved}.

\begin{thebibliography}{9}
\bibitem{Chevalley1954}
C.~Chevalley, \emph{The Algebraic Theory of Spinors}, Columbia University Press, New York, 1954.

\bibitem{Lounesto2001}
P.~Lounesto, \emph{Clifford Algebras and Spinors}, 2nd ed., London Mathematical Society Lecture Note Series, vol.~286, Cambridge University Press, Cambridge, 2001.

\bibitem{Porteous1995}
I.~R.~Porteous, \emph{Clifford Algebras and the Classical Groups}, Cambridge Studies in Advanced Mathematics, vol.~50, Cambridge University Press, Cambridge, 1995.
\end{thebibliography}

\end{document}
